\section{A common interval of backward survival}
\label{sec:survival}

We now turn the restricted pair-correlation theorem into a common backward interval for a large family of simple zeros. The local estimate depends on the entire complementary divisor through a positive inverse-square sum.

\begin{lemma}[One zero on a complex time disk]\label{lem:one-zero-survival}
Let $f$ be real entire of order at most one with real zeros, and let $x$ be a simple zero. Define
\[
 \ell_x=\frac{f''(x)}{2f'(x)},\qquad
 S_x=\sum_{f(y)=0,\,y\ne x}(x-y)^{-2}.
\]
For every $V>0$ with $S_x\le V$, the zero has a simple holomorphic continuation throughout $|a|\le2^{-16}/V$ under $e^{aD^2}f$. It obeys
\[
 z_x(-a)=x-2a\ell_x+w_x(a),\qquad
 |w_x(a)|<\frac1{32\sqrt V},
\]
and is real for real $a$ in the interval.
\end{lemma}
\begin{proof}
The centered product of Lemma~\ref{lem:centered-product} gives
$f(x+h)=f'(x)h e^{\ell_xh}R(h)$ and
$|R(h)-1|\le8V|h|^2e^{2V|h|^2}$. Conjugating the heat operator yields
\[
 \frac{f_{-a}(x-2a\ell_x+w)}{f'(x)e^{\ell_xw-a\ell_x^2}}
 =G_a(w)=\mathbb E[(w+\sqrt{2a}Z)R(w+\sqrt{2a}Z)],
\]
where $Z$ is standard Gaussian. As before, Gaussian domination and evenness in the square root give holomorphic dependence on $a$.
Set $\kappa=2^{-5}$, $c_0=2^{-16}$ and $r=\kappa/\sqrt V$. For $|w|=r$, write $v=\sqrt Vw$ and $\sigma=\sqrt{2aV}$. Then
\[
 \sqrt V|G_a(w)-w|
 \le8\mathbb E[|v+\sigma Z|^3e^{2|v+\sigma Z|^2}].
\]
On $|a|\le c_0/V$, use $|v|=\kappa$, $|\sigma|\le\sqrt{2c_0}$,
$(u+v)^3\le4(u^3+v^3)$, and
$2|v+\sigma Z|^2\le4\kappa^2+8c_0Z^2$.
The right side is at most
\[
 32e^{4\kappa^2}\bigl(
 \kappa^3\mathbb Ee^{8c_0Z^2}
 +(2c_0)^{3/2}\mathbb E|Z|^3e^{8c_0Z^2}\bigr).
\]
The elementary Gaussian integrals
$\mathbb Ee^{Z^2/4}=\sqrt2<2$ and
$\mathbb E|Z|^3e^{Z^2/4}=16/\sqrt{2\pi}<8$, together with
$e^{4\kappa^2}<2$, bound this by
$128\kappa^3+512(2c_0)^{3/2}<\kappa/4$.
Rouch\'e's theorem gives exactly one zero in $|w|<r$, counting multiplicity. It is consequently simple. Its implicit sections agree on overlaps by uniqueness and extend through the boundary of the parameter disk because the estimate is strict. Conjugation fixes the unique section on the real parameter interval. This proves all assertions.
\end{proof}

Let $N_H(X)$ count the positive zeros of $H_0$ up to $X$, with multiplicity, and put $M_T=N_H(2T)-N_H(T)$. The zero-location input and the Riemann--von Mangoldt formula give
$N_H(X)\sim X\log X/(4\pi)$.

Put
\[
 c_{\rm MT}=\frac32-\frac1{\sqrt2}\cot\frac1{\sqrt2}
            =0.6725007\ldots.
\]
This is the Montgomery--Taylor simple-zero proportion, whose smooth
test formulation is given in \cite[Lemma~3.2, pp.~13--15]{lamzouri-simple2026}.

\begin{theorem}[Common survival at the Montgomery--Taylor proportion]\label{thm:two-thirds-survival}
For every $0<\eta<c_{\rm MT}$ there exist $c_\eta>0$ and $T_\eta$ such that, for all $T\ge T_\eta$, at least $(c_{\rm MT}-\eta)M_T$ distinct simple zeros of $H_0$ in $(T,2T)$ have branches which remain distinct, real, simple, and in $(T,2T)$ throughout
\[
                  -\frac{c_\eta}{(\log T)^2}\le t\le0.
\]
Consequently, for every positive $a(T)$ with $a(T)(\log T)^2\to0$, the lower limiting proportion of initial centers having those properties on $[-a(T),0]$ is at least $c_{\rm MT}$.
\end{theorem}
\begin{proof}
We first state the precise correlation input. For each nonnegative even Schwartz function $g$ with $g(0)=1$ and Fourier support strictly inside $(-1,1)$, the degree-one case of Rudnick--Sarnak \cite[Theorems~1.1--1.2, pp.~4--5]{rudnick-sarnak}, using their positive-zero normalization and their multiplicity convention on pp.~1--3, gives
\[
 P(Y):=\sum_{\substack{0<x_j,x_k\le Y\\j\ne k}}
       g(F(x_j/2)-F(x_k/2))
       =I_gN_H(Y)+o(N_H(Y)),
\]
where
\[
 F(u)=\frac{u\log u}{2\pi},\qquad
 I_g=\int_{\mathbb R}g(u)
           \left(1-\left(\frac{\sin\pi u}{\pi u}\right)^2\right)du.
\]
The sum is over distinct indices, so repeated ordinates contribute.
There are admissible $g$ with $I_g<1-c_{\rm MT}+\eta/16$.
We verify the precise normalization of the Montgomery--Taylor test
used here. On $[-1/2,1/2]$ put
\[
 f_0(v)=\frac{\cos(\sqrt2v)}{\sqrt2\sin(1/\sqrt2)},
\]
and extend it by zero. It is nonnegative and has integral one.
For $g_0(u)=|\int f_0(v)e^{-2\pi iuv}\,dv|^2$, Plancherel and
Fourier inversion of the squared sinc give
\[
 I_{g_0}=\int f_0(v)^2\,dv-1+
       \iint |v-w|f_0(v)f_0(w)\,dv\,dw.
\]
The function $f_0(v)+\int|v-w|f_0(w)dw$ has second derivative
$f_0''+2f_0=0$ in the interior and is even, so it is constant.
At $v=1/2$ its value is
$1/2+\cot(1/\sqrt2)/\sqrt2=2-c_{\rm MT}$.
Integrating this identity against $f_0$ gives
$I_{g_0}=1-c_{\rm MT}$.
Approximate $f_0$ in $L^1\cap L^2$ by even nonnegative smooth
functions supported strictly inside $(-1/2,1/2)$, each normalized
to have integral one. Their squared Fourier moduli are nonnegative
Schwartz functions, have value one at zero and Fourier support
strictly inside $(-1,1)$. The $L^2$ convergence of the transforms
implies $L^1$ convergence of their squared moduli, hence convergence
of $I_g$, since the sinc weight is bounded. This proves the asserted
admissible approximation and the exact energy conversion.
Fix one sufficiently good smooth approximation. Choose $0<\delta<1$ with $g(u)\ge1-\eta/16$ for $|u|\le\delta$, and put $h=\delta/\log T$.

Call a zero occurrence in $(T,2T)$ isolated when every other occurrence has distance at least $h$. Such an occurrence is simple. Since
\[
 \frac{d}{dx}F(x/2)=\frac{\log(x/2)+1}{4\pi}\le\log T
 \quad(T-1\le x\le2T+1)
\]
for large $T$, every nonisolated occurrence has an ordered witness pair contributing at least $1-\eta/16$. Select one witness per first index; these ordered pairs are distinct. They all occur in
$P(2T+1)-P(T-1)$, whose asymptotic is $I_gM_T+o(M_T)$.
Thus the proportion of nonisolated occurrences is at most
\[
 \frac{I_g+o(1)}{1-\eta/16}
 <1-c_{\rm MT}+\frac\eta4
\]
eventually. At least $(c_{\rm MT}-\eta/4)M_T$ isolated centers remain.

We next control $S_x$ on most of these centers. For any fixed complementary occurrence $y\in(T/2,3T)$, the isolated centers on either side are mutually $h$-separated and, after omitting coincidence, are all at distance at least $h$ from $y$. Hence
\[
 \sum_{x\text{ isolated},\,x\ne y}(x-y)^{-2}
 \le2h^{-2}\sum_{k\ge1}k^{-2}<4h^{-2}.
\]
Summing this inequality over $y$ preserves every complementary multiplicity. Since $N_H(3T)\le4M_T$ eventually, the middle region contributes at most $16M_Th^{-2}$. The zero-count bound and partial summation show uniformly for $x\in(T,2T)$ that the zeros outside $(T/2,3T)$, including the negative zeros, contribute $O(\log T/T)$ to $S_x$. Therefore
\[
                 \sum_{x\text{ isolated}}S_x\le20M_Th^{-2}
\]
for large $T$. Put $B=80/(\eta\delta^2)$. Markov's inequality discards at most $\eta M_T/4$ centers with $S_x>B\log^2T$.

The bound $N_H(u+1)-N_H(u)=O(\log(u+3))$ follows directly from the positive Hadamard sum
\[
 \sum_\gamma\frac{3/2}{(3/2)^2+(v-\gamma)^2}
 =\Re\frac{\xi'}\xi(2+iv)=O(\log(v+3)).
\]
Here the last estimate uses the logarithmic derivative of the gamma factor and the absolutely convergent Euler series at real part $2$. Each occurrence with $|v-\gamma|\le1$ contributes at least $6/13$. Consequently the endpoint strips of width $\log T$ contain $O(\log^2T)=o(M_T)$ zeros. Remove them as well; at least $(c_{\rm MT}-\eta)M_T$ centers remain for large $T$.

Apply Lemma~\ref{lem:one-zero-survival} with the common value $V=B\log^2T$. To control the drift, the even canonical product gives
\[
 \ell_x=\frac1{2x}+
 \sum_{y>0,\,y\ne x}
       \left(\frac1{x-y}+\frac1{x+y}\right).
\]
The paired terms with $y<x/2$ or $y>2x$ sum absolutely to $O(\log T)$. In the middle region the terms $1/(x+y)$ have the same bound. The terms at distance below $1$ are at most
$(O(\log T)S_x)^{1/2}=O_B((\log T)^{3/2})$ by Cauchy's inequality. Summing the remaining unit-distance shells gives $O(\log^2T)$. Thus $|\ell_x|=O_B(\log^2T)$ uniformly on the selected set. The displacement on $|a|\le2^{-16}/V$ is $O_B(1)$, so the endpoint margin keeps each branch in $(T,2T)$.

Distinct selected branches cannot meet. A meeting at a simple zero would identify their local implicit sections, and the identity theorem on the connected common time disk would identify their different initial centers. Take $c_\eta=2^{-16}/B$.
Finally, for each fixed $\eta$, the condition $a(T)\log^2T\to0$ places $[-a(T),0]$ within the proved interval for all sufficiently large $T$. Taking the lower limit and then $\eta\downarrow0$ gives the final assertion.
\end{proof}
